\documentclass{article}
\usepackage[utf8]{inputenc}
\usepackage[T1]{fontenc}
\usepackage{geometry}
\geometry{
    a4paper,
    left=2.5cm,
    right=2.5cm,
    top=2.5cm,
    bottom=2.5cm
}
\usepackage{graphicx}
\usepackage{float}
\usepackage{hyperref}
\usepackage[english]{babel}
\usepackage{amsmath}
\usepackage{amssymb}
\usepackage{array}
\usepackage{booktabs}
\usepackage{caption}
\usepackage{subcaption}
\usepackage{enumitem}
\usepackage{fancyhdr}
\usepackage{titlesec}
\usepackage{url}
\usepackage{textcomp}
\usepackage{xcolor}
\usepackage{listings}
\usepackage{multirow}
\usepackage{longtable}
\usepackage{tabularx}
\usepackage{times}

% Set up page style
\pagestyle{fancy}
\fancyhf{}
\rfoot{\thepage}
\renewcommand{\headrulewidth}{0pt}

% Title formatting
\titleformat{\chapter}[display]{\normalfont\bfseries}{\chaptertitle\thechapter}{0pt}{\Huge}
\titlespacing*{\chapter}{0pt}{50pt}{40pt}

% Document info
\title{Template Document}
\author{Academic Support Agent}
\date{\today}

\begin{document}
\maketitle

%%%%% N o alterar esta quatro linhas abaixo %%%%%%%
    \begin{frame}{} 
        \titlebox
        \vspace{0.7cm}
    \begin{columns} 
    
    
    %%%%%% inicio coluna 1 %%%%%%%%
    \begin{columnbox} % divis o em coluna
        \begin{block}{\Large Introdu o}
            \justifying
            \blindtext[3] % comando para preenchimento de texto gerado automaticamente
        \end{block}
        
        \blindtext[1]
        
        \begin{figure}
            \centering
            \includegraphics[scale=0.1]{logo/bottom.png}
            \caption{Simbolo SEFiM}
            \label{fig:my_label}
        \end{figure}
        
        \begin{block}{\Large Revis o Bibliogr fica}
            \begin{table}[]
                \centering
                \begin{tabular}{c|c}
                     lorem & ipsun \\
                     dolor & sit
                \end{tabular}
                \caption{Tabela de Lorem Ipsun}
                \label{tab:my_label}
            \end{table}
        \end{block}
        
    \end{columnbox} % finaliza o de divis o em coluna
    
     %%%%%% inicio coluna 2 %%%%%%%%
    \begin{columnbox} % divis o em coluna
        \begin{block}{\Large Resultados}
            \justifying
            \blindtext[4] % comando para preenchimento de texto gerado automaticamente
        \end{block} 
        \blindtext[2]
    \end{columnbox} % finaliza o em coluna
    
     %%%%%% inicio coluna 3 %%%%%%%%
    \begin{columnbox} % divis o em coluna
        \begin{block}{\Large Conclus o}
            \justifying
            \blindtext[3] % comando para preenchimento de texto gerado automaticamente
        \end{block}
        \blindtext[1]
        \begin{block}{\Large Refer ncias}
            \begin{thebibliography}{9}
            % altere os dados abaixo para as suas referencias
                \bibitem{latexcompanion} 
                Michel Goossens, Frank Mittelbach, and Alexander Samarin. 
                \textit{The \LaTeX\ Companion}. 
                Addison-Wesley, Reading, Massachusetts, 1993.
                 
                \bibitem{einstein} 
                Albert Einstein. 
                \textit{Zur Elektrodynamik bewegter K{\"o}rper}. (German) 
                [\textit{On the electrodynamics of moving bodies}]. 
                Annalen der Physik, 322(10):891--921, 1905.
                 
                \bibitem{knuthwebsite} 
                Knuth: Computers and Typesetting,
                \\\textit{http://www-cs-faculty.stanford.edu/\~{}uno/abcde.html}
            \end{thebibliography}
        \end{block}
        
        \begin{block}{\Large Agradecimentos}
            \justifying
            % altere os dados abaixo para os seus agradecimentos
            \begin{itemize}
                \item Primeiro item 
                \item Segundo item
                \item Terceiro item
                \item Quarto item 
            \end{itemize}
        \end{block}
        
    \end{columnbox}

    \end{columns} % finaliza o em coluna
    \end{frame}
\end{document}